\documentclass[10pt,a4paper]{amsart}
\usepackage{amsmath,amssymb,amsthm,graphicx}
\usepackage{verbatim}
\newenvironment{reference}{\comment}{\endcomment}
%\newenvironment{reference}{}{}
\newenvironment{slogan}{\comment}{\endcomment}
\newenvironment{history}{\comment}{\endcomment}

% For commutative diagrams we use Xy-pic
\usepackage[all]{xy}

% We use 2cell for 2-commutative diagrams.
\xyoption{2cell}
% Theorem environments.
%
\theoremstyle{plain}
\newtheorem{theorem}[subsection]{Theorem}
\newtheorem{proposition}[subsection]{Proposition}
\newtheorem{lemma}[subsection]{Lemma}

\theoremstyle{definition}
\newtheorem{definition}[subsection]{Definition}
\newtheorem{example}[subsection]{Example}
\newtheorem{exercise}[subsection]{Exercise}
\newtheorem{situation}[subsection]{Situation}

\theoremstyle{remark}
\newtheorem{remark}[subsection]{Remark}
\newtheorem{remarks}[subsection]{Remarks}

\numberwithin{equation}{subsection}

% Macros
%
\def\lim{\mathop{\mathrm{lim}}\nolimits}
\def\colim{\mathop{\mathrm{colim}}\nolimits}
\def\Spec{\mathop{\mathrm{Spec}}}
\def\Hom{\mathop{\mathrm{Hom}}\nolimits}
\def\Ext{\mathop{\mathrm{Ext}}\nolimits}
\def\SheafHom{\mathop{\mathcal{H}\!\mathit{om}}\nolimits}
\def\SheafExt{\mathop{\mathcal{E}\!\mathit{xt}}\nolimits}
\def\Sch{\mathit{Sch}}
\def\Mor{\mathop{\mathrm{Mor}}\nolimits}
\def\Ob{\mathop{\mathrm{Ob}}\nolimits}
\def\Sh{\mathop{\mathit{Sh}}\nolimits}
\def\NL{\mathop{N\!L}\nolimits}
\def\CH{\mathop{\mathrm{CH}}\nolimits}
\def\proetale{{pro\text{-}\acute{e}tale}}
\def\etale{{\acute{e}tale}}
\def\QCoh{\mathit{QCoh}}
\def\Ker{\mathop{\mathrm{Ker}}}
\def\Im{\mathop{\mathrm{Im}}}
\def\Coker{\mathop{\mathrm{Coker}}}
\def\Coim{\mathop{\mathrm{Coim}}}

% Boxtimes
%
\DeclareMathSymbol{\boxtimes}{\mathbin}{AMSa}{"02}

%
% Macros for moduli stacks/spaces
%
\def\QCohstack{\mathcal{QC}\!\mathit{oh}}
\def\Cohstack{\mathcal{C}\!\mathit{oh}}
\def\Spacesstack{\mathcal{S}\!\mathit{paces}}
\def\Quotfunctor{\mathrm{Quot}}
\def\Hilbfunctor{\mathrm{Hilb}}
\def\Curvesstack{\mathcal{C}\!\mathit{urves}}
\def\Polarizedstack{\mathcal{P}\!\mathit{olarized}}
\def\Complexesstack{\mathcal{C}\!\mathit{omplexes}}
% \Pic is the operator that assigns to X its picard group, usage \Pic(X)
% \Picardstack_{X/B} denotes the Picard stack of X over B
% \Picardfunctor_{X/B} denotes the Picard functor of X over B
\def\Pic{\mathop{\mathrm{Pic}}\nolimits}
\def\Picardstack{\mathcal{P}\!\mathit{ic}}
\def\Picardfunctor{\mathrm{Pic}}
\def\Deformationcategory{\mathcal{D}\!\mathit{ef}}
\usepackage{cleveref}
\crefname{originallemma}{original lemma}{original lemmas}
\Crefname{originallemma}{Original lemma}{Original lemmas}
\makeatletter
\newlabel{lemma-categorically-compact-QCoh}{{0GPE}{source}{}{sourcepointer-1}{}}
\newlabel{lemma-categorically-compact-QCoh@cref}{{[originallemma][0][]0GPE}{[1][0][]source}{}{}{}}
\makeatother
\crefname{originalreference}{original reference}{original references}
\Crefname{originalreference}{Original reference}{Original references}
\makeatletter
\newlabel{modules-lemma-finite-presentation-stalk-free}{{0B8J}{source}{}{sourcepointer-2}{}}
\newlabel{modules-lemma-finite-presentation-stalk-free@cref}{{[originalreference][0][]0B8J}{[1][0][]source}{}{}{}}
\makeatother
\crefname{originalreference}{original reference}{original references}
\Crefname{originalreference}{Original reference}{Original references}
\makeatletter
\newlabel{topology-lemma-quasi-compact-closed-point}{{005E}{source}{}{sourcepointer-3}{}}
\newlabel{topology-lemma-quasi-compact-closed-point@cref}{{[originalreference][0][]005E}{[1][0][]source}{}{}{}}
\makeatother
\crefname{originalreference}{original reference}{original references}
\Crefname{originalreference}{Original reference}{Original references}
\makeatletter
\newlabel{schemes-section-quasi-coherent}{{01LA}{source}{}{sourcepointer-4}{}}
\newlabel{schemes-section-quasi-coherent@cref}{{[originalreference][0][]01LA}{[1][0][]source}{}{}{}}
\makeatother
\crefname{originalreference}{original reference}{original references}
\Crefname{originalreference}{Original reference}{Original references}
\makeatletter
\newlabel{coherent-lemma-quasi-compact-h1-zero-covering}{{01XF}{source}{}{sourcepointer-5}{}}
\newlabel{coherent-lemma-quasi-compact-h1-zero-covering@cref}{{[originalreference][0][]01XF}{[1][0][]source}{}{}{}}
\makeatother
\title{If QCoh(X) is equivalent to a module category, X is affine}
\author{The Stacks Project Authors}
\date{stacks-project:functors.tex:local-snapshot-20261012; GNU FDL 1.2 or later}
\begin{document}
\maketitle
\noindent\textit{Editorial note.} All delivered mathematical text is the author's own, copied verbatim from the pinned Stacks Project chapter file functors.tex; no statement or proof step was rewritten, reordered or removed.\par\medskip
\noindent\textit{Editorial note.} Span recomputation by label pairing (LANE\_BRIEF 10.2): in this chapter the author's proof environment follows the statement environment as a standalone block; the spec covers the complete environment pair (statement span, proof span) as recorded in the pool candidate.\par\medskip
\noindent\textit{Editorial note.} Occupation scan: combined claim\_key + statement\_anchor + generic-line scan over all registration files (active and withdrawn) found 0 hits for functors.tex (fu-occupied.json, 2026-10-14); consistent with the duty manager's six-check dispatch.\par\medskip
\ifdefined\hypertarget\hypertarget{sourcepointer-1}{}\fi
\noindent\textit{Source reference.} Original lemma 0GPE: sources/stacks-chow/tags.txt, source lines 20824--20824. Cross-chapter standard prerequisite cited by the author inside the counted span; it is used as a pointer and its own proof is not part of this item.\par\medskip
\ifdefined\hypertarget\hypertarget{sourcepointer-2}{}\fi
\noindent\textit{Source reference.} Original reference 0B8J: sources/stacks-chow/tags.txt, source lines 14142--14142. Cross-chapter standard prerequisite cited by the author inside the counted span; it is used as a pointer and its own proof is not part of this item.\par\medskip
\ifdefined\hypertarget\hypertarget{sourcepointer-3}{}\fi
\noindent\textit{Source reference.} Original reference 005E: sources/stacks-chow/tags.txt, source lines 206--206. Cross-chapter standard prerequisite cited by the author inside the counted span; it is used as a pointer and its own proof is not part of this item.\par\medskip
\ifdefined\hypertarget\hypertarget{sourcepointer-4}{}\fi
\noindent\textit{Source reference.} Original reference 01LA: sources/stacks-chow/tags.txt, source lines 2079--2079. Cross-chapter standard prerequisite cited by the author inside the counted span; it is used as a pointer and its own proof is not part of this item.\par\medskip
\ifdefined\hypertarget\hypertarget{sourcepointer-5}{}\fi
\noindent\textit{Source reference.} Original reference 01XF: sources/stacks-chow/tags.txt, source lines 2516--2516. Cross-chapter standard prerequisite cited by the author inside the counted span; it is used as a pointer and its own proof is not part of this item.\par\medskip
\begin{lemma}
\label{lemma-characterize-affine}
Let $X$ be a quasi-compact and quasi-separated scheme.
If $\QCoh(\mathcal{O}_X)$ is equivalent to the category
of modules over a ring, then $X$ is affine.
\end{lemma}

\begin{proof}
Say $F : \text{Mod}_R \to \QCoh(\mathcal{O}_X)$ is an equivalence.
Then $\mathcal{F} = F(R)$ has the following properties:
\begin{enumerate}
\item it is a finitely presented $\mathcal{O}_X$-module
(Lemma \ref{lemma-categorically-compact-QCoh}),
\item $\Hom_X(\mathcal{F}, -)$ is exact,
\item $\Hom_X(\mathcal{F}, \mathcal{F})$ is a commutative ring,
\item every object of $\QCoh(\mathcal{O}_X)$ is a quotient of
a direct sum of copies of $\mathcal{F}$.
\end{enumerate}
Let $x \in X$ be a closed point. Consider the surjection
$$
\mathcal{O}_X \to i_*\kappa(x)
$$
where the target is the pushforward of $\kappa(x)$ by the
inclusion morphism $i : x \to X$. We have
$$
\Hom_X(\mathcal{F}, i_*\kappa(x)) =
\Hom_{\mathcal{O}_{X, x}}(\mathcal{F}_x, \kappa(x))
$$
This first by (4) implies that $\mathcal{F}_x$ is nonzero.
From (2) we deduce that every map $\mathcal{F}_x \to \kappa(x)$
lifts to a map $\mathcal{F}_x \to \mathcal{O}_{X, x}$ (as it
even lifts to a global map $\mathcal{F} \to \mathcal{O}_X$).
Since $\mathcal{F}_x$ is a finite $\mathcal{O}_{X, x}$-module,
this implies that $\mathcal{F}_x$ is a (nonzero)
finite free $\mathcal{O}_{X, x}$-module.
Then since $\mathcal{F}$ is of finite presentation, this implies
that $\mathcal{F}$ is finite free of positive rank
in an open neighbourhood of $x$
(Modules, Lemma \ref{modules-lemma-finite-presentation-stalk-free}).
Since every closed subset of $X$ contains a closed point 
(Topology, Lemma \ref{topology-lemma-quasi-compact-closed-point})
this implies that $\mathcal{F}$ is finite locally free of positive rank.
Similarly, the map
$$
\Hom_X(\mathcal{F}, \mathcal{F}) \to
\Hom_X(\mathcal{F}, i_*i^*\mathcal{F}) =
\Hom_{\kappa(x)}(\mathcal{F}_x/\mathfrak m_x \mathcal{F}_x,
\mathcal{F}_x/\mathfrak m_x \mathcal{F}_x)
$$
is surjective. By property (3) we conclude that the rank $\mathcal{F}_x$
must be $1$. Hence $\mathcal{F}$ is an invertible $\mathcal{O}_X$-module.
But then we conclude that the functor
$$
\mathcal{H}
\longmapsto
\Gamma(X, \mathcal{H}) = \Hom_X(\mathcal{O}_X, \mathcal{H}) =
\Hom_X(\mathcal{F}, \mathcal{H} \otimes_{\mathcal{O}_X} \mathcal{F})
$$
on $\QCoh(\mathcal{O}_X)$ is exact too. This implies that the first
$\Ext$ group
$$
\Ext^1_{\QCoh(\mathcal{O}_X)}(\mathcal{O}_X, \mathcal{H}) = 0
$$
computed in the abelian category $\QCoh(\mathcal{O}_X)$ vanishes for all
$\mathcal{H}$ in $\QCoh(\mathcal{O}_X)$. However, since
$\QCoh(\mathcal{O}_X) \subset \textit{Mod}(\mathcal{O}_X)$
is closed under extensions
(Schemes, Section \ref{schemes-section-quasi-coherent})
we see that $\Ext^1$ between quasi-coherent modules
computed in $\QCoh(\mathcal{O}_X)$ is the same as computed
in $\textit{Mod}(\mathcal{O}_X)$. Hence we conclude that
$$
H^1(X, \mathcal{H}) =
\Ext^1_{\textit{Mod}(\mathcal{O}_X)}(\mathcal{O}_X, \mathcal{H}) = 0
$$
for all $\mathcal{H}$ in $\QCoh(\mathcal{O}_X)$.
This implies that $X$ is affine for example by
Cohomology of Schemes, Lemma
\ref{coherent-lemma-quasi-compact-h1-zero-covering}.
\end{proof}

\end{document}
